\section{Continuous paths and their finite certificates}\label{sec:continuity}
A continuous curve can visit a face infinitely often. A real endpoint on a mirror can be arbitrarily hard to distinguish from nearby endpoints on either side. Both issues admit finite estimates.

\paragraph{A short bridge between nearby samples.}
Use a supporting affine form $s_G\ge0$ for every face $G$, zero on $G$, scaled so its linear part has norm at most one. Choose a positive rational $\sigma$ smaller than every strictly positive vertex slack $s_G(v)$. This is finite algebraic data. For example, bound a positive slack below by rationalizing or isolating the finitely many fixed radicals in its expression. Let $b$ be an upper bound on face diameters (two suffices here). If $p\in F$, $q\in G$ and $|p-q|\le\eta<\sigma$, then $s_G(p)\le\eta$. Express $p$ as a convex combination of vertices of $F$. The total weight outside $G$ is at most $\eta/\sigma$. There must be a common vertex. Move all outside weight to a fixed such vertex. The result $r\in F\cap G$ satisfies $|p-r|\le b\eta/\sigma$. The chords $pr,rq$ form a surface bridge of length at most $(1+2b/\sigma)\eta$. A face-center weight can be distributed equally among the vertices.

For these four models one can take $b=2$, $\sigma=1/4$ throughout. For a tetrahedral or octahedral face take its outward sign normal divided by two, whose norm is $\sqrt3/2$; the positive vertex slack is $1/\sqrt2$. For the icosahedron use the unscaled coordinates $w=2p$ and the normal $n=(2\phi-2,0,-2\phi)$. The support value is $2+2\phi$ in $w$-coordinates, and its positive slacks there are $4,4\phi,4+4\phi$. The physical form with linear part $n/4$ has slacks divided by eight and $\|n/4\|=\sqrt{12}/4$. For face zero of the dodecahedron in the unscaled coordinates $w=2p/\phi$, take $n=(2-2\phi,2\phi-4,0)$. The support value is $2$, and the positive slacks there are $8-4\phi,4\phi-4,4$. In physical coordinates the linear part $n/4$ has squared norm $(28-16\phi)/16<1$ and the slacks are multiplied by $\phi/8$. The smallest is $(\phi-1)/2>1/4$. Substitution of the vertex lists verifies these finite tables. Face transitivity transfers them to the other faces.

\paragraph{Compressing a path with arbitrarily many crossings.}
Write $f$ for the number of faces. Sample a continuous curve at a mesh on which consecutive samples differ by at most $\delta$. Approximate each sample to ambient tolerance $\delta$ by a tagged surface point, retaining exact given endpoints when available. Consecutive tagged points differ by at most $3\delta$. Start at the first tag. Jump to the \emph{last} index with the same face tag, join the two tagged points by a chord, and bridge to the next tagged point by the preceding construction. Repeat from that next index. No finished face tag reappears, so at most $f$ chords and $f$ bridges are used. The sum of the chords is at most the sum of the original ambient sample increments plus $2f\delta$, by the triangle inequality on disjoint index intervals. The resulting surface polygonal path has length at most
\[
  \sum_i|\gamma(t_{i+1})-\gamma(t_i)|
    + f\,[2+3(1+2b/\sigma)]\delta,\qquad 3\delta<\sigma .
\]
The explicit choices give the error $53f\delta$. Taking $\delta=2^{-(k+11)}$ makes it less than $2^{-k}$ for $f\le20$. The sampled sum is at most the length of the original curve. Since $\delta$ is arbitrary, finite-path lower bounds hold for all rectifiable continuous curves. This proves also the asserted pathwise bound for a scalar function that is $1$-Lipschitz on convex pieces and agrees at their boundaries: split each face chord at its finitely many piece boundaries and sum the estimates. Recover exact endpoints by a short bridge if they were approximated. In the first-order interpretation a modulus supplies the mesh, and the tags belong to the input approximants.

\paragraph{Computing the intrinsic distance at prescribed precision.}
For exact tagged algebraic endpoints enumerate all lists of distinct faces whose first and last faces contain the given endpoints and whose consecutive intersections are nonempty. Include a single common face. For every shared edge allow crossings at the $2^m+1$ equally spaced grid points, and for a shared vertex allow that vertex alone. Join successive choices by face chords. Minimize the sum of their lengths rounded \emph{up} to multiples of $2^{-m}$. Call the result $U_m$. These candidates are genuine surface paths. Conversely, shorten any finite path by removing repeated faces via chords, then round its at most $f-1$ crossings to the grid on the shared intersection. This changes its length by at most $2(f-1)2^{-m}$ for unit edges. Upward length rounding contributes at most $f2^{-m}$. The two path arguments give
\[
 d\le U_m\le d+3f\,2^{-m}.
\]
Changing the endpoints by surface distances $\epsilon_x,\epsilon_y$ changes $d$ by at most $\epsilon_x+\epsilon_y$. For ambient approximations use the short-bridge estimate. A prescribed dyadic accuracy thus determines one finite search. The number of candidates is at most $\sum_{r=1}^f f^r(2^m+1)^{r-1}$, including illegal lists rejected by the checker. This algorithm defines the distance as a regular Cauchy name. At every extremal pair the proofs give a specific path attaining the distance.

\section{Arithmetizing the proof in iPRA}
Use intuitionistic first-order natural-number arithmetic with primitive recursive function symbols, their defining equations and quantifier-free induction, as in ProveIt. A statement about a real object denotes a family of bounds on finite inputs and demanded precisions. A path supplied by primitive recursive sample and modulus functions gives a schema, one instance for each choice of functions. The geometric interpretation applies externally to ordinary continuous paths with their moduli. First-order iPRA has no variable ranging over all real-valued functions.

\paragraph{Data types and exact comparisons.}
Encode signed integers by differences of naturals, rationals by a signed numerator and positive denominator, and finite lists by primitive recursive pairing and length. A solid is a fixed record of vertex coordinates, oriented face cycles and flag triples. A mesh point has a flag index and nonnegative barycentric numerators $a_0,a_1,a_2$ with positive sum $A$, meaning $\sum_i a_i C_i/A$. Weights with sum $2^m$ give a mesh: round down two weights and give the remainder to the third. The ambient error is at most $2b\,2^{-m}$. Finer meshes represent general endpoints without deciding their exact face. Interior waypoints and seams may use exact algebraic coordinates. A path code records its face list and endpoints of each segment. Test incidence with plane equations and polygon halfspaces, and equality of the two physical copies of each seam point.

Use a fixed tower of rational arithmetic and positive square roots for the model constants, for example $\sqrt2,\sqrt3,\sqrt5,\sqrt{10+2\sqrt5}$. Mirror reflections have the formula $p\mapsto p-2(n\cdot p)n/(n\cdot n)$. For pentagons the smaller ring $\mathbb Z[z]$ suffices for developed vertices and hinges, with $\operatorname{Re}(z)=(\sqrt5-1)/4$ and $\operatorname{Im}(z)=\sqrt{10+2\sqrt5}/4$. To test the sign of $a+b\sqrt r$ at a quadratic level with $r>0$, compare the signs of $a,b$ first. For opposite signs compare $a^2$ and $b^2r$ at the previous level. Represent general quotients by numerator and nonzero denominator, carrying its sign and cross-multiplying. The recursion depth of these comparisons is fixed. For a squared segment length $s$, its dyadic square-root floor is the largest integer $j$ with $j^2\le2^{2m}s$. For face chords $j\le2^{m+1}$. For a developed strip take a bound from its fixed number of faces and their diameters. The ceiling uses one further unit exactly when needed. A long sum of segment lengths needs only dyadic sums and explicit errors.

\paragraph{Object-language assertions.}
For a fixed solid $S$, a claimed pair $w,w^*$ and the positive radical $D_S$ on page one, a purely numerical specification at exact tagged endpoints is
\begin{align*}
 &\forall p,q,k\quad \mathrm{Point}_S(p)\land\mathrm{Point}_S(q)
     \ \Longrightarrow\ U_{k+c}(p,q)\le D_S+2^{-k},\\
 &\forall P,k\quad \mathrm{Path}_S(P,w,w^*)
     \ \Longrightarrow\ \sum_{e\in P}2^{-k}\lceil 2^k|e|\rceil\ge D_S .
\end{align*}
One may take $c=6$ simultaneously for the four solids, since $2^6\ge3f$. Clear positive denominators and square after checking signs to compare to $D_S$. The lower test accepts coded algebraic facewise paths of arbitrary finite length. A fixed witness strip and the mesh estimate give equality of distance names. For the upper test the inequalities give routes whose best length is at most $D_S$; rounding their crossings gives the displayed finite bound. One can compare the two resulting dyadic costs, so route selection on a real-order tie costs at most its rounding tolerance. For Cauchy input add the two explicitly controlled endpoint errors and shift the precision index. Optimal coordinates are fixed algebraic points.

\paragraph{Constructive equality cases.}
For the icosahedron and dodecahedron the gaps bound $1-\lambda_1\mu_1$. Since both weights lie in $[0,1]$, each endpoint's weight outside its vertex corner is bounded by that deficit. For the octahedron, if $d^2$ is within $\epsilon$ of $3$, each sum $\sum_{i<j}a_i a_j$ is at most $\epsilon$. Choose a largest coordinate on a finite approximant, at least $1/3$. The other coordinates sum to at most $3\epsilon$. In all three cases both endpoints are close to vertices, and nonantipodal vertices have the explicit strict upper bounds given above.

For the tetrahedron, divide each nearest-vertex quadrilateral by the diagonal joining its lattice vertex to the large triangle center $h'$. If the barycentric weight on $h'$ is $t$, convexity gives $\delta^2\le (1-t)+4t/3$. Thus $\delta^2\ge4/3-\epsilon$ places the point within $6\epsilon$ of $h'$, as the triangle diameter is at most two. Apply this to both $u+v$ and $u-v$. Their two center residues are within $6\epsilon$ of signed translates of $h$. Adding and subtracting identifies one endpoint within $6\epsilon$ of a lattice vertex and the other within $6\epsilon$ of its opposite-face center: equal signs give $v$ near a lattice vertex and $u$ near that vertex plus or minus $h$ modulo $2L$, and opposite signs reverse the roles. Folding is nonexpansive. This leaves a finite list of isolated candidate pairs in every case. Choose one precision at which neighborhoods of distinct candidates are disjoint, and keep the two approximation errors smaller than a quarter of that separation. All finer approximations must select the same unordered pair. Thus equality for names gives membership in the finite list without an exact real sign test on a mirror.

\paragraph{Permitted inductions.}
A finite fold, gallery, list minimum or bounded sign check has a primitive recursive trace and Boolean invariants. In path compression the index increases and the used-face bit mask grows. In a strip search the unused-face count decreases. A universal over indices below the loop counter is a bounded Boolean conjunction. Apply quantifier-free induction to the trace and its invariant. Rational triangle and Jensen inequalities follow after clearing denominators from nonnegative squares: Cauchy--Schwarz uses $\sum_{i<j}(u_i v_j-u_j v_i)^2$, and the variance identity subtracts $\sum_{i<j}\lambda_i\lambda_j\|u_i-u_j\|^2$. Inequalities involving square roots of segment norms may be checked on dyadic enclosures to arbitrary precision, allowing their rounding errors per segment. The compression argument bounds this count independently of the number of original samples. No compactness axiom selects an extremizer.

\section{Lean verification and formalization roadmap}
The supplied Lean sources were checked with Lean 4.32.0. Two kinds of theorem are included.
\begin{itemize}
\item \texttt{SmallCertificates.lean} checks the tetrahedral hole and the octahedral corner and permutation tables. \texttt{Icosahedron.lean} checks the two-strip matrix and the calibration tile measurements.
\item \texttt{DodecaLower.lean} enumerates the fifty strips by kernel reduction, checks the minimum multiplicity, the attaining hinge identities, and the four-edge neighborhood. \texttt{DodecaUpper.lean} constructs the four-face strip, checks all 27 intersections and the transposed gap matrix. These four files use core Lean and exact integer computation.
\item \texttt{IntrinsicIPRA.lean} and \texttt{IcosahedronIPRA.lean} construct actual derivations in ProveIt's embedded theory, of type $\Gamma\vdash_{\rm PRAAx}\varphi$. Their \texttt{octMinimum} and \texttt{icoMinimum} prove the strip-minimum bounds for arbitrary natural barycentric masses after clearing denominators. Other theorems derive a two-weight variance identity and the coordinate formula and multiplicativity of the cyclotomic squared norm. The Lean metatheory of this provider uses quotient soundness, propositional extensionality and classical choice. The object proofs use the iPRA axioms.
\item \texttt{TetrahedronIPRA.lean} derives the upper bound on a Voronoi triangle and the lower bound for every signed lattice translate at the hole. \texttt{DodecahedronIPRA.lean} derives the weighted gap bound, including the term $\frac3{10}(1-\lambda_1\mu_1)$, for the matrix as a function of a rational parameter at least one. Substitution of arbitrarily accurate rational approximants to $\phi$ gives the stated numerical interpretation.
\end{itemize}
The remaining Lean work is the geometric interpretation of these finite objects and assembly of the full object-language diameter theorems. A dependency order follows.
\begin{enumerate}
\item Build PR codes for ordered algebraic arithmetic, support forms, binary interval costs and tagged cells. Use the provider's \texttt{Core/Pairing}, \texttt{Core/Order}, \texttt{Core/Division} and bounded search modules. Prove the sign and rounding contracts one quadratic level at a time. Give realizations of the face tables, with both labelings of each shared edge. The exact dodecahedral support-plane checks are in \texttt{dodeca\_model.py}.
\item Define the predicate on facewise paths specified above. Derive chord replacement, gluing, the short bridge and rounding. Prove the precision bound on the enumerated minimum by pushing each crossing onto an edge mesh and bounding its two adjoining length changes. This connects a sample/modulus instance to the lower test.
\item Encode chamber wall reflections, the finite gallery fold, and interpolation on hinges. Invariants are the unresolved wall signs, membership of each crossing in its hinge, and the running bound on path length. Use norm identities and variance to compile the average bounds. For the tetrahedron encode lifts in a bounded fundamental parallelogram and enumerate shifts satisfying the quadratic bound.
\item Translate each fixed Boolean certificate as a closed PR term. Derive what a successful return implies about the geometry using Steps 1--3. ProveIt's \texttt{eval\_closed} in \texttt{Core/Evaluation} derives the corresponding numeral computation in the object theory. Combine with the attaining route, strict gaps and isolation argument to close the universal statements. The sampling theorem and this last assembly are provided here as mathematical proofs.
\end{enumerate}
\paragraph{Sources.}
The two attached cuboid reports supply the problem and the interpretation of surface distance. For the farthest-point map at arbitrary starting points, see J. Rouyer, \emph{Antipodes sur un t\'etraèdre r\'egulier}, J. Geom. 77 (2003), 152--170, and R. E. Schwartz's papers on the \href{https://arxiv.org/abs/2004.11852}{regular octahedron} and the \href{https://arxiv.org/abs/2104.02567}{regular dodecahedron}. The diameter calculations and certificates used here appear in the preceding sections. The object-theory implementation targets the \href{https://github.com/VladimirReshetnikov/ProveIt/tree/09fe17194da14cdd100947737952e44100bb71b8/Logic/PrimitiveRecursiveArithmetic}{ProveIt iPRA development}, revision \texttt{09fe17194d}.
